\section{Introduction}
\label{sec:intro}
% ------------------------------------------------------------
% Sumber: §13.1-13.3, §13.8.3, simple_notes/2 & 3.
% Alur: (1) hook Sim2Real gap, (2) dua literatur berhenti satu langkah
% sebelum jembatan, (3) pertanyaan kita, (4) tiga lapis, (5) kontribusi.
% Dipangkas 2026-09-13 (restrukturisasi front matter, Agent A): <=600 kata,
% 5 kontribusi -> 3 (lihat RESTRUCTURE_PLAN.md).
% ------------------------------------------------------------

% --- Paragraf 1: hook — simulasi survei LLM itu dipakai orang, tapi outputnya homogen
LLMs are increasingly deployed as simulated survey respondents and human
simulators \citep{subpop2025,odyssim2026}. Yet behavioral evidence shows
their outputs are excessively uniform: simulated users are ``overly
agreeable'' and ``stylistically uniform'' \citep{sim2real2026}, opinion
predictions barely differentiate demographic groups, and general model
capability does not predict simulation fidelity \citep{sim2real2026}. This
evidence is \emph{behavioral}---measured from the outside. Whether the
failure lives in what the model \emph{knows} or in what it \emph{uses} has
remained open, because the two research communities each hold only half
of the required instrumentation.

% --- Paragraf 2: dua literatur, masing-masing setengah alat (analogi MRI/kuesioner)
Interpretability work localizes demographic and persona information inside
models---down to individual attention heads \citep{llmopinions2026,
kimevans2025}---but validates against the model's own behavior (probe
accuracy, steering success), never against external population ground
truth. LLM-survey-simulation work has exactly that ground truth (real
per-group response distributions) but treats the model as a black box.
The experiment that requires both instruments at once---whether the
internal arrangement of groups matches how those groups actually
answer---sits in the gap between them and had not been run.

% --- Paragraf 3: pertanyaan riset + instrumen
We hold both instruments at once. For 169 intersectional demographic
cells (e.g., \textit{Asian $\times$ Hindu}, \textit{age 18--29 $\times$
Democrat}) we compute two representational dissimilarity matrices---one
from real Pew ATP response distributions (Wasserstein distances between
groups' answers), one from the model's internal representations at a
given location---and correlate them (RSA; \citealp{kriegeskorte2008}).
Applied at 1{,}089 locations $\times$ 4 prompt operationalizations, this
yields a \emph{fidelity map}: where inside the model does the arrangement
of demographic groups mirror how the groups actually differ in opinion?

% --- Paragraf 4: tiga lapis (tesis utama)
Our results split the question ``can LLMs simulate populations?'' into
three separable properties (we reserve ``layer'' for transformer layers):
\begin{enumerate}
  \item \textbf{Readable} (established by prior work): per-group opinion
        distributions can be decoded from internal activations
        \citep{llmopinions2026}.
  \item \textbf{Faithfully arranged} (this work): at specific locations,
        including a single attention head, \starhead, the \emph{inter-group
        geometry} mirrors real opinion structure (selection-corrected
        $\rrho$ up to $0.63$), which per-group readability neither implies
        nor requires.
  \item \textbf{Actually used} (this work): the faithfully arranged
        information barely flows into the model's answers---full identity
        swaps in the prompt move predictions by ${<}2\%$ of their error,
        and not directionally toward the target group's truth. A
        strong-instrument intervention does find a real causal pathway,
        but where it is clearest is \emph{not} where the map is most
        faithful (Sec.~\ref{sec:sweep}).
\end{enumerate}

% --- Paragraf 5: kontribusi eksplisit (5 butir lama -> 3)
\paragraph{Contributions.}
\begin{itemize}
  \item A \emph{fidelity map} of demographic identity in an LLM---to our
        knowledge the first to score internal inter-group geometry against
        real survey ground truth: RDM-vs-RDM comparison across layers,
        attention heads, FFN outputs, and prompt variants, with
        winner's-curse-corrected statistics (max-statistic permutation,
        held-out and split-half validation)---and identifies \starhead{}
        as a general demographic-geometry head, significant across all
        six attribute types as a \emph{fixed} location, with fidelity
        heterogeneous by type: political/socio-economic structure is
        encoded at $\rrho$ up to $0.63$, while racial$\times$religious
        structure is weak and prompt-fragile everywhere.
  \item A dissociation of causal use from representational fidelity.
        Replacing the entire prompt identity moves the model's predicted
        distribution by under $2\%$ of its prediction error in aggregate
        (Sec.~\ref{sec:gap}), and a causal sweep over all 32 layers in all
        six types---with cluster-robust, selection-corrected
        inference---shows that the clearest causal pathway sits in a
        low-fidelity type while the most faithful type shows no
        correction-surviving single-layer pathway (Sec.~\ref{sec:sweep}).
  \item A direct measurement of what a faithful read-out buys, with a
        practical reframing. A probe on \starhead{} alone (128 dimensions)
        is $21$--$31\%$ closer to survey truth than the model's own
        answers, but recovers almost no per-question group ordering: the
        two measure different resolutions, faithful \emph{on average} but
        not sharp enough for any single question (Sec.~\ref{sec:probe}).
        If faithful structure already exists internally, surface
        fine-tuning \citep{subpop2025} may re-teach what the model already
        encodes---a candidate mechanism for its known transfer
        failure---and reading the location directly is the cheaper
        alternative, though our probe results show it is still not a
        substitute for survey data.
\end{itemize}
